\section{El lenguaje no es la esencia de la inteligencia, sino un salto cualitativo}

Esta sección parte del concepto que más se presta a confusión. En la entrevista, Wang He (王鹤) define la inteligencia de un modo cercano a «la capacidad de responder al entorno según la situación». Esta definición no pone el lenguaje en el centro de forma deliberada, porque muchos seres vivos carecen de un lenguaje de orden superior y aun así muestran inteligencia en su entorno; la visión, el tacto, el sonido y la propiocepción pueden ser los sensor con los que un agente inteligente percibe el mundo. La importancia del lenguaje reside en que permite a los seres humanos comprimir la experiencia, difundir el conocimiento, formar conceptos abstractos y aprender de manera colaborativa; por eso es un salto cualitativo, y no la única fuente de la inteligencia.

Este juicio incide directamente en la orientación de la inteligencia corporizada. Si se equipara la inteligencia con el lenguaje, se tiende a pensar que los LLM ya resolvieron la mayor parte de los problemas; si la inteligencia se entiende como la capacidad de responder al entorno, se advierte que el robot debe volver a la percepción, la acción y la retroalimentación. Un modelo de lenguaje puede aportar sentido común y descomponer tareas, pero el robot tiene que actuar en el mundo físico, donde un fallo altera el entorno y tiene costos; esto exige otro tipo de datos y otra forma de evaluación.

\begin{figure}[H]
\centering
\includegraphics[width=0.94\textwidth]{language-intelligence-jump.png}
\caption{El lenguaje no es la esencia de la inteligencia: es un salto cualitativo de la inteligencia humana, y no la única vía de acceso a la inteligencia. Diagrama conceptual propio, elaborado a partir de la conversación de 00:05:58--00:25:08.}
\end{figure}

\begin{knowledgebox}{Lectura de la figura: el lenguaje potencia la inteligencia, pero no sustituye al cuerpo}
A la izquierda, intelligence pone el acento en el entorno, la respuesta y la retroalimentación; a la derecha, language lo pone en la compresión, la difusión y la colaboración. Al leer esta figura, no hay que ver ambos lados como opuestos, sino como una relación jerárquica: el cuerpo y la percepción proporcionan el contacto con el entorno, y el lenguaje permite que la experiencia se transmita entre individuos. El problema de la inteligencia corporizada es que no basta con que el lado del lenguaje sea fuerte; el lado físico también debe cerrar el ciclo.
\end{knowledgebox}

\subsection{La visión es un sensor potente, no una inteligencia completa}

Esta subsección distingue con más detalle entre visión e inteligencia. La visión es un sensor muy potente, que apareció antes que muchas otras modalidades perceptivas en la evolución de la inteligencia de los animales superiores; sin embargo, el reconocimiento puramente visual suele ser passive perception, es decir, dada una imagen, determinar su categoría, segmentar regiones o estimar posiciones. Una vez concluido el reconocimiento, el entorno no cambia porque uno diga «esto es un gato», ni proporciona retroalimentación para el siguiente paso. Esa es la diferencia entre las tareas de visión basadas en internet y las tareas de inteligencia corporizada.

Wang He señala que el campo de la visión por computadora acumuló durante años una enorme capacidad gracias a tareas como ImageNet, el reconocimiento facial y la segmentación semántica, pero estas tareas estaban impulsadas en su mayoría por datos de internet y anotación manual. Cuando la competencia en esas tareas llegó a cierto grado de saturación, los investigadores empezaron a buscar el siguiente paso: hacer que la visión fuera activa, que el agent emprendiera acciones, que las acciones modificaran las observaciones y que las nuevas observaciones guiaran el siguiente paso. De ahí surge el Perception-Action Loop.

\begin{warningbox}{No hay que equiparar la capacidad de reconocimiento con la capacidad corporizada}
Poder reconocer una taza, una silla y una habitación no equivale a poder desplazarse junto a la silla, tomar la taza, abrir la puerta de un gabinete u ordenar un anaquel. La inteligencia corporizada requiere acción y retroalimentación del entorno, no solo asignar etiquetas a imágenes estáticas.
\end{warningbox}

\subsection{Resumen de la sección}

El lenguaje produjo un salto en la inteligencia humana, pero la inteligencia corporizada no puede depender solo del lenguaje. El robot necesita enlazar la visión, la acción y la retroalimentación del entorno en un ciclo cerrado; este es también el motor principal del giro de la investigación en visión por computadora hacia la inteligencia corporizada.
